\section{Training data and design of experiments}
\label{sec:doe}

The choice of training data, sampling design and fidelity strategy is an
integral part of the surrogate model. A design of experiments (DoE)
specifies which input configurations are evaluated, and, in adaptive
settings, in what order and at which model fidelity, to build or refine
the surrogate. For MES in particular, that design space and sampling
plan which the optimizer will explore, covers multiple energy carrier
inputs.

This section reviews the four DoE choices that recur most often in our
analyzed body of literature -- data sources, static space-filling
designs, adaptive sampling, and multi-fidelity training -- and discusses
the associated limitations.

\begin{table*}[!t]
\centering
\scriptsize
\begin{tabularx}{\textwidth}{@{}L{0.18\textwidth} L{0.30\textwidth} L{0.20\textwidth} L{0.26@\textwidth} Strategy & Description & Pairs naturally with@{}}
\toprule
& Typical task (refs.)\
Latin hypercube sampling & Stratified random fill of the design
hyper-cube when no informative prior exists & Kriging, RBF, GP &
Stochastic dispatch emulation \cite{Hu2021671}\
Sobol / quasi-MC & Low-discrepancy or quadrature points paired with
polynomial-chaos expansions & PCE, sparse grid, GP & Chance-constrained
OPF and stochastic economic dispatch
\cite{Muhlpfordt20174490,Safta20172535,Koirala20242284,Xu20212696}\
Factorial / fractional FD & Discrete grid or structured experimental
design over a small factor set & RSM, GPR, low-order PCE & Geothermal
design under uncertainty and iterative plant/controller optimization
\cite{Schulte2020,Deodhar2018}\
Adaptive sampling & Iteratively adds truth evaluations via infill rules
(error, trust region) & Kriging, RBF, GP & Combined dispatch, OPF,
and expensive system-design search \cite{Lin2023,Deng2020831}\
Active learning (BO) & Acquisition-driven query (EI / LCB) selecting the
next high-fidelity evaluation & GP (preferred), deep ensembles & Unit
commitment, building/HVAC calibration
\cite{Nikolaidis2021,Chakrabarty2021}\
Multi-fidelity training & Combines cheap low-fidelity sweeps with sparse
high-fidelity evaluations & Co-kriging, multi-fidelity GNN, MF GP & OPF,
geothermal design, battery prognostics
\cite{Taghizadeh2024,Chen2025244,Valladares2022}\
Transfer learning & Fine-tunes a pretrained surrogate on a small sample
from a new context & NN, GBDT, GP & Battery state-of-charge across
ambient conditions \cite{Chen202514792}\
Counterfactual / synthetic & Labels from repeated high-fidelity model
evaluation at prescribed inputs & NN / L2O proxies, RSM, GP & SCED/SCUC,
OPF, stability-constrained dispatch proxies
\cite{Wu2022231,Chen20244723,Chen2022,Liu20223726,Liang20246508}\
Historical operational data & Past measurements; serial correlation and
feature leakage if not handled & NN, GBDT, GP residuals & Load, wind and
renewable uncertainty modelling and downstream planning
\cite{Hu2021671,Safta20172535,Koirala20242284,Chakrabarty2021}\
\end{tabularx}
\end{table*}

\subsection{Data sources and workflow context}
\label{sec:doe:sources}

Based on the literature, two recurrent training-data sources for the
DoE can be identified -- synthetic and historical -- together with a
third recurring integration mode in which surrogate modeling is embedded
in a hybrid workflow that combines data-driven surrogates with
simulation- or model-based system components. These three cases are
described in the following.

\subsubsection{Synthetic data}

Optimizer-synthetic data is generated by repeatedly evaluating a
high-fidelity model. Because those outputs come from the same
high-fidelity model that the surrogate is meant to replace, the
surrogate is trained to return the same objective values or other model
results that the optimizer would obtain from that model at each
candidate design. In power-system studies, Liu et al. sample the search
space and evaluate each candidate dispatch with full small-signal
stability analysis before learning an explicit optimal power flow
constraint \cite{Liu20223726}. Xu et al. use correlated uncertainty
scenarios and build a PCR surface from repeated optimal power flow runs
\cite{Xu20212696}. Chen et al. repeatedly evaluate a power-flow simulator at
varying load and renewable generation to build a training set, then fit
quantile predictors for chance-constrained optimal power flow without
explicit network parameters \cite{Chen2023657}. Wu et al., Chen et al. and
Liang et al. train proxies from large batches of economic dispatch for
market clearing under varying load, renewable output and cost data
\cite{Wu2022231,Chen2022,Liang20246508}. For buildings, Thrampoulidis et
al. run multiple simulations for building retrofit and insulation in
Zurich and train a surrogate from the resulting costs and system
selections \cite{Thrampoulidis2021}. Aghaei Pour et al. use the same idea
where the surrogate first narrows the search, but the detailed
simulation is rerun only for configurations that the planner chooses to
check during interactive design \cite{Aghaei Pour2022303}. At district
scale, Sánchez-Zabala et al. use repeated EnergyPlus simulations to
train building metamodels for district energy-management optimization
\cite{Sánchez-Zabala2024}, and Chaudhry et al.
evaluate a high-fidelity fifth-generation DH and cooling model over
uncertain markets and demands to train a surrogate inside robust design
optimization \cite{Chaudhry2024}. Related MES studies of Liu et al. use
Latin-hypercube sampling and refine a kriging model with adaptive infill
\cite{Liu201914569}. Hirvonen et al. replace repeated solar-community
simulations with a neural surrogate model during MOO sizing
\cite{Hirvonen2017323} and Gao et al. analyze battery, electrolyzer and
solid oxide fuel cell sizes in MATLAB--TRNSYS co-simulation \cite{Gao2025}.

\subsubsection{Historical data}

historical data are the prevalent source for forecasting surrogates,
because they describe how an energy system actually operated in the past
rather than how a simulation model would respond to a new design. For
MES, Shi and Teh forecast several load types in a regional MES by
splitting each measured history into quick fluctuations, slower
variations and a long-term trend, predicting each part separately and
adding the results back together \cite{Shi2024__2}.\cite{Dong2024}. Koschwitz et
al. predict monthly heating and cooling loads for a German district from
building consumption data using support-vector and recurrent networks
\cite{Koschwitz2018134}, while Bianchi et al. extend autoregressive models
of DH load in Italy with convolutional and Gaussian-process predictors
for horizons up to 48Â~h \cite{Bianchi2020244}. Also forecasting of day-ahead
campus electricity load \cite{Madhukumar20228891}, short-term wind speed and
wind power \cite{Tian2020,Hanifi2024}, and photovoltaic output \cite{Li2019}
has been conducted.

\subsubsection{Hybrid workflows}

Hybrid workflows combine data-driven surrogates with simulation- or
model-based system components in the wider optimization or control
loop. In these studies, the distinguishing feature is not necessarily a
mixed training-data source, but rather the coupling of learned
approximations to a broader model-based workflow. Sánchez-Zabala et
al. train building metamodels from repeated EnergyPlus simulations and
couple them to optimization algorithms in a district energy-management
testbed \cite{Sánchez-Zabala2024}. Reich et al. fit approximations from selected
recent measurements in a local DH network and embed them in dynamic
simulation for predictive operation \cite{Reich2020}. Jalving et al. embed
market surrogates trained from a library of annual power-market
simulations into MES design optimization \cite{Jalving2023}. Dong et al.
fit a PCR model from historical MES operation and use it in day-ahead
dispatch \cite{Dong2023}. Yin et al. combine operating records with a
surrogate model to size electricity and hydrogen storage in a MES
\cite{Yin2024}. This category therefore captures workflow coupling rather
than only the provenance of the surrogate training data.

\subsection{Static designs: LHS, Sobol, factorial}
\label{sec:doe:static}

As the optimum is unknown prior to optimization, the first training set
must explore the full search space.Static experimental designs fix that
set before any surrogate is fitted. They differ in how points are spread
through the input space and in what surrogate class they are meant to
support. Three static design classes can be identified based on our
literature analysis: Latin hypercube sampling (LHS), quasi-Monte Carlo
and sparse-grid collocation, and factorial or response-surface designs.

\subsubsection{Latin hypercube sampling}

LHS divides each input dimension into equal intervals and places exactly
one sample in each interval, with random permutations across dimensions.
Compared with unstructured random sampling, LHS gives more even coverage
of samples. Unlike factorial designs it does not require discrete factor
levels ((a fixed set of values per input, such as low/medium/high,
rather than any point in a continuous range) and keeps the number of
runs moderate. It is therefore widely used to initialize surrogates such
as kriging or Gaussian processes before optimization.

For a MES, Liu et al.Â~use Latin hypercube samples to fit a kriging model
of annual cost \cite{Liu201914569}. Seele et al.Â~solve MILP only at LHS to
calibrate surrogates of total annualized cost for distributed
energy-system planning under uncertainty, \cite{Seele20232480}. Summ et
al.Â~employ LHS from a thermo-hydraulic collector model to train a
surrogate model before MOO for solar DH \cite{Summ2025}. In single-carrier
dispatch surrogate modeling, Hu et al.Â~apply LHS in a reduced
wind-uncertainty space to train a Gaussian-process model of stochastic
economic dispatch \cite{Hu2021671}.

\subsubsection{Quasi-Monte Carlo, sparse grids, and polynomial-chaos collocation}

When optimization parameter oder data such as renewable generation,
loads, or prices are uncertain, a stochastic optimizer often needs
expected costs or violation probabilities. The Monte Carlo approach
draws random values of the uncertain inputs, solve the energy model once
for each draw, and average the results. That becomes expensive when each
solve is a full dispatch or network calculation. Quasi-Monte Carlo
methods replace purely random draws with low-discrepancy point sets
(Sobol and Halton sequences are common examples) that cover the
uncertainty space more evenly. In our analyzed body of literature,
however, this appears as sparse grids or quadrature/collocation nodes
used to build a PCE: each node is one deterministic simulator run, and
the PCE coefficients summarize how outputs respond to uncertain inputs.
While LHS spreads design inputs across a box to train a surrogate of
performance indicators before optimization, sparse-grid and PCE designs
instead employ uncertainty scenarios so the optimizer can estimate
expected costs, output spread, or violation probabilities inside a
stochastic formulation without thousands of repeated solves.

For MES, Dong et al.Â~use a data-driven sparse PCE to approximate
day-ahead dispatch responses under chance constraints, learning from
historical data so that iterative verification does not require repeated
full-order solves \cite{Dong2023}. Jiang et al.Â~apply PCE to approximate the
security-region boundary of MES, replacing repeated nonlinear solves
during boundary search \cite{Jiang2026620}. Gaitanis et al.Â~quantify
capacity uncertainty in HRES with CHP for buildings using PCE after
design optimization \cite{Gaitanis2026}. Safta et al.Â~propagate wind
uncertainty in stochastic economic dispatch with PCE and sparse
quadrature instead of Monte Carlo sampling \cite{Safta20172535}, Lu et
al.Â~use nested sparse-grid collocation under wind and photovoltaic
uncertainty in stochastic economic dispatch \cite{Lu201991827}, Muhlpfordt
et al.Â~formulate chance-constrained AC optimal power flow with PCE
\cite{Muhlpfordt20194806}, Engelmann et al.Â~combine PCE with distributed
stochastic optimal power flow \cite{Engelmann20186188}, and Xu et al.Â~use a
polynomial-chaos response surface with dependent uncertain inputs for
chance-constrained AC optimal power flow \cite{Xu20212696}.

\subsubsection{Factorial and response-surface designs}

Factorial and response-surface designs are well-suited for problems with
only a few adjustable settings, for example plant size, temperature
levels, or cooling technology, and where it matters how those settings
combine: does a larger unit still help when the operating temperature is
already high? A full factorial design runs the simulator at every
combination of discrete factor levels (a fixed set of values per input,
such as low/medium/high, rather than any point in a continuous range).
Fractional factorial, Box--Behnken, and related layouts reduce the run
count while still estimating main effects and selected interactions. The
fitted surface is typically a low-order polynomial metamodel. Unlike
LHS, the inputs are chosen from those predefined levels rather than
spread continuously through the search space; unlike sparse-grid or PCE
designs, the goal is to explore a small set of design parameters rather
than to propagate input uncertainty through a stochastic formulation.

Salahinezhad et al.Â~run a CCHP plant model over the operating period for
each candidate design, record energy, cost and emission indicators, and
fit a response surface to those outputs before optimizing trigeneration
layout with solar and hydrogen subsystems \cite{Salahinezhad2025}. Liu et
al.Â~use Box--Behnken designs and response-surface to analyze how
equipment capacity, price and energy coefficients affect annual cost and
carbon emissions in a MES model \cite{Liu2025__4}. Roy et al.Â~employ
response surface methodology to MOO of a biomass CHP plant \cite{Roy2023}
and Ahmad et al.Â~use a Box--Behnken design over catalyst amount,
electrode voltage and electrolysis time to model green-hydrogen
production from wastewater electrolysis \cite{Ahmad202451}. Cao et al.Â~apply
experimental design with response-surface regression to select microgrid
schemes under stochastic short-term simulation \cite{Cao2017} and El Mestari
et al.Â~use a full factorial $2^2$ design to tune genetic-algorithm
crossover and mutation coefficients for wind-farm layout
\cite{El Mestari2025}

\subsection{Adaptive sampling and active learning}
\label{sec:doe:adaptive}

Static designs from the previous subsection fix all simulation points
before the surrogate is trained and used. Adaptive methods add further
points during the optimization process, e.g. where the surrogate is
inaccurate or near a constraint limit. Based on our literature review,
we distinguish between infill rules that plug into surrogate-assisted
search, and acquisition-driven queries in Bayesian optimization.

\subsubsection{Infill rules}

Infill rules choose the next expensive model run from the current
surrogate state. For a MES, Liu et al.Â~start from LHS and then add
kriging infill points using minimum surrogate value and
mean-squared-error criteria until the annual-cost model is accurate
enough \cite{Liu201914569}. Lin et al.Â~update kriging surrogates with a new
infilling rule while searching economic dispatch on large power systems
\cite{Lin2023}. Deng et al.Â~deploy kriging so that most candidate optimal
power flow evaluations use the surrogate and only promising points
trigger a full AC power flow solve \cite{Deng2020831}. In building
energy-system design, Aghaei Pour et al.Â~use stated preference MOO
search to select which candidate configurations receive a full
simulation for surrogate update \cite{Aghaei Pour2022303}. Granacher et
al.Â~predict whether each candidate layout of an industrial utility
system is Pareto-optimal, and run the detailed process model next where
that prediction is most uncertain \cite{Granacher20221573}.

\begin{table*}[!t]
\centering
\scriptsize
\begin{tabularx}{\textwidth}{@{}L{0.04\textwidth} L{0.16\textwidth} L{0.16\textwidth} L{0.18\textwidth} L{0.20\textwidth} L{0.18@\textwidth} \# & Pattern name & Surrogate@{}}
\toprule
role & Solver layer & Safety hooks & Typical pitfalls\
P1 & Replace expensive subroutine & Direct functional substitute for an
inner optimizer or simulator & Inner loop / subproblem &
Feasibility-restoration layer; constraint-aware embedding; periodic
truth-model re-anchoring
\cite{Liu20223726,Subedi2026157,Liu20242534} & Silent constraint
violation; OOD operating regimes; surrogate-induced bias of the host
objective\
P2 & Accelerate search & Acquisition / pre-screening of expensive
candidates & Outer loop / metaheuristic / BO & Truth-model verification
of every selected candidate; calibrated posterior; trust-region updates
\cite{Nikolaidis2021,Schulte2020,Aghaei Pour2022303} &
Mis-calibrated posterior, uninformative acquisition, wasted
truth-budget\
P3 & Warm start / solution proxy & Predicts (near-)optimal decisions
from problem data & Pre-solve / branch-and-bound seeding & Downstream
solver as repair step; projection onto feasible set; ensembling for
variance reduction
\cite{Liang20246508,Wu2022231,Chen20244723,Chen2022} &
Constant-time output is infeasible; warm start dominated by repair cost\
P4 & Decomposition support & Approximates recourse / pricing / value
function & Benders / column generation / DP layer & Periodic exact cuts;
bound tracking on optimality gap; restart on cut staleness
\cite{Wu2024391,Liu20242534,Li2023__2,Zhang2025__4} & Surrogate cuts
dominate but are not certified; non-convergence of master--subproblem
loop\
P5 & Uncertainty handling & Approximates chance / robust / DRO mapping &
Reformulation layer & Calibrated tails (PIT, reliability diagrams);
sample-based audit of chance constraints; PCE moment matching
\cite{Muhlpfordt20194806,Safta20172535,Hu2020,Hu2021671,Volodina2022}
& Tail under-coverage; Gaussian-only assumptions; correlation collapse
across operating regimes\
\end{tabularx}
\end{table*}

\subsubsection{Bayesian optimization and active learning}
\label{sec:doe:bo_al}

Bayesian optimization (BO) fits a Gaussian process (or related
probabilistic surrogate). Here, it is assessed via an acquisition
function where to simulate next: it favours inputs that look strong on
the surrogate and inputs where the surrogate is still uncertain, and the
highest-scoring candidate receives the next high-fidelity run. Griffith
et al.Â~apply this framework to techno-economic optimization of a MES,
reporting fewer expensive evaluations \cite{Griffith2025}. Chakrabarty et
al.Â~calibrate 17 coupled building and HVAC parameters with scalable BO,
dividing the search space as the cost surface is learned
\cite{Chakrabarty2021}. Uhrich et al.Â~use active learning to pick training
runs for electricity-market metamodels, reporting large runtime
reductions with limited training data \cite{Uhrich2026211}. Matta et al.Â~use
active learning with k-nearest-neighbour and Gaussian-process surrogates
to find near-optimal solutions in islanded microgrids after only a few
full evaluations \cite{Matta2026}. Nikolaidis et al.Â~use GP-BO for economic
dispatch under volatile renewable production \cite{Nikolaidis2021}.

\subsection{Multi-fidelity training and transfer learning}
\label{sec:doe:multifidelity}

Many energy-system optimization studies have models at different
computational cost and accuracy levels, for example, a detailed network
simulationm, reduced-order models, simple analytic approximations, etc.
Multi-fidelity training mixes these data were many cheap runs plus a few
expensive ones are combined, so the surrogate learns where the cheap
model is trustworthy and where the expensive model must be used.
Transfer learning addresses a related problem: the expensive model is
unchanged, but the operating context shifts, for example temperature or
market conditions, and a surrogate trained in one setting is updated
with a small number of runs in the new one.

Chung and Wang study a nuclear--renewable hybrid plant that also feeds
an industrial process: a detailed Simulink model supplies high-fidelity
training data and a fast forward calculation supplies low-fidelity data,
and a co-kriging surrogate trained on both predicts control settings
across timescales \cite{Chung2023}. Taghizadeh et al.Â~train a GNN on
multiple low-fidelity and fewer high-fidelity runs, thus, reducing
training cost while keeping accuracy for uncertain points
\cite{Taghizadeh2024}. Khayambashi et al.Â~apply an enhanced multi-fidelity
GNN surrogate to chance-constrained optimal power flow
\cite{Khayambashi202453}, and Xu et al.Â~propagate large sample sets through
cheap polynomial-chaos surrogates and use the high-fidelity power-system
model only to correct rare high-risk events \cite{Xu20203127}.

While multi-fidelity training mixes cheap and expensive versions of the
same model, transfer learning keeps one model and reuses a surrogate
already trained in one context when a similar new context appears (such
as changed market conditions, temperaetures, etc.). The goal is to avoid
repeating the initial surrogate training, hence, most of the training
data is retained, and only a few new simulations in the new conditions
are needed to adapt it. Chen et al.Â~illustrate this for battery
state-of-charge estimation: a recurrent network with a Gaussian-process
layer is first trained at one ambient temperature and then adapted
quickly when the temperature changes \cite{Chen202514792}.

\subsection{Limitations specific to energy data}
\label{sec:doe:pitfalls}

Surrogate models fail quietly when training data do not match how the
system is operated or validated. The following major limitations from
the analyzed literature can be identified.

First, Du et al.Â~report that variability in training datasets is often
overlooked when fitting surrogates for DH with thermal storage, even
though operating strategy strongly affects robustness under peak demand
\cite{Du2025}. Second, Faraji et al.Â~use load and weather forecasts in
day-ahead microgrid scheduling \cite{Faraji2020157284} and a review of
learning-assisted power-system decision making describes exogenous
forecasts as inputs to optimization problems \cite{Lim2025}. Third,
Sánchez-Zabala et al.Â~note that metamodels generalize only as far as the
quality and availability of their training simulations allow
\cite{Sánchez-Zabala2024}. Extreme heat can dominate cooling requirements
\cite{Morakinyo201973}, and probabilistic risk studies keep the full
power-system model for rare high-risk tail events that cheap surrogates
mishandle \cite{Xu20203127}.
